% ============================================================
%  ch07 —— §5.1–5.2（原文 p.26–p.31）
% ============================================================
\routenote{p.26--p.31}

\section{Kloosterman 层的上同调性质}

本节设 $G$ 是 $k$ 上分裂几乎单群，视作 $\PP^1$ 上的常群概形。回忆对每个根 $\alpha\in\Phi$，$U_\alpha$ 是 $G$ 中的根群。对每个根 $\alpha$ 固定同构 $u_\alpha:\mathbb{G}_a\xrightarrow{\sim}U_\alpha$。

\medskip\noindent\textbf{记号.} 对 $k$ 上概形 $X$ 与 $X$ 上 $\mathbb{Q}_\ell$-层复形 $F$，$\chi_c(X,F)$ 与 $H^*(X,F)$ 指 $F$ 拉回到 $X\otimes_k\bar k$ 后的 Euler 特征与上同调。

\subsection{Euler 特征、Swan 导子与上同调刚性}

有两个维数较小的 $V\in\Rep(\ckG)$ 的一致选择。其一是 $V=\mathrm{Ad}=\ckg$，即伴随表示；其二是 $V=V_{\theta^\vee}$，即非零权由 $\ckg$ 的短根组成的表示。称 $V_{\theta^\vee}$ 为 $\ckG$ 的\emph{拟极小表示}（当 $\ckG$ 单链时它与伴随表示重合）。关于 $V_{\theta^\vee}$ 的基本事实汇总在附录 §10。记 $G$ 的长（相应地短）单根个数为 $r_\ell(G)$（相应地 $r_s(G)$），称之为 $G$ 的长秩（相应地短秩）。这些数的等价描述见引理 10.1。

记 $\Kl^{\theta^\vee}_{\ckG}(\varphi,\chi)$（相应地 $\Kl^{\mathrm{Ad}}_{\ckG}(\varphi,\chi)$）为与 $\Kl_{\ckG}(\varphi,\chi)$ 及 $\ckG$ 的拟极小表示（相应地伴随表示）相伴的局部系统。

\begin{mainthm}[4]
\begin{enumerate}[label=(\arabic*), leftmargin=1.9em, itemsep=.25em]
\item $-\chi_c(\PP^1\setminus\{0,\infty\},\Kl^{\theta^\vee}_{\ckG}(\varphi,\chi))$ 等于 $G$ 的长单根个数。
\item 设当 $G$ 非单链时 $\operatorname{char}(k)$ 对 $G$ 良好，则 $-\chi_c(\PP^1\setminus\{0,\infty\},\Kl^{\mathrm{Ad}}_{\ckG}(\varphi,\chi))$ 等于 $G$ 的秩。
\end{enumerate}
\end{mainthm}

由 Grothendieck--Ogg--Shafarevich 公式，得到

\begin{corollary}
\begin{enumerate}[label=(\arabic*), leftmargin=1.9em, itemsep=.2em]
\item $\Swan_\infty(\Kl^{\theta^\vee}_{\ckG}(\varphi,\chi))=r_s(\ckG)$。
\item 在与定理 4(2) 相同的 $\operatorname{char}(k)$ 假设下，$\Swan_\infty(\Kl^{\mathrm{Ad}}_{\ckG}(\varphi,\chi))=r(\ckG)$。
\end{enumerate}
\end{corollary}

以下各小节将致力于证明定理 4。先推出一些推论。

\begin{lemma}
设主幂幺元作用在 $V$ 上无平凡 Jordan 块，则 $H^0(\PP^1\setminus\{0,\infty\},\Kl^V_{\ckG}(\varphi))=0$。
\end{lemma}

\begin{proof}
设不然，则 $\Kl^V_{\ckG}(\varphi)$ 以常层为子局部系统。按构造 $\Kl^V_{\ckG}(\varphi)$ 是纯的。由 [3, Th.~5.3.8]，在 $\PP^1\setminus\{0,\infty\}\otimes_k\bar k$ 上 $\Kl^V_{\ckG}(\varphi)$ 是单 perverse 层之直和。故该常子层必是直和项。

另一方面，按假设，$V$ 不含任何使 $\PP^1$ 中 $\{0\}$ 处的驯单值（由定理 1，即主幂幺元）平凡作用的直和项。这导致矛盾。
\end{proof}

设 $i_0:\{0\}\hookrightarrow\PP^1$，$i_\infty:\{\infty\}\hookrightarrow\PP^1$，$j:\PP^1\setminus\{0,\infty\}\hookrightarrow\PP^1$ 为包含。对 $\PP^1\setminus\{0,\infty\}$ 上任意局部系统 $L$，滥用记号，以 $j_{!*}L$ 表示 $j_{!*}(L[1])[-1]$。

如定理 2 之前的讨论那样，局部 Galois 群 $\operatorname{Gal}(K^{\mathrm{sep}}_\infty/K_\infty)$ 与 $\operatorname{Gal}(K^{\mathrm{sep}}_0/K_0)$（从而惯性群 $I_0$ 与 $I_\infty$）作用在 $\Kl^V_{\ckG}(\varphi)$ 的相应几何茎上，定义 $V$ 上在共轭意义下的表示。

\begin{proposition}
在与定理 4(2) 相同的 $\operatorname{char}(k)$ 假设下，有
\begin{enumerate}[label=(\arabic*), leftmargin=1.9em, itemsep=.2em]
\item （上同调刚性）$H^*(\PP^1_k,j_{!*}\Kl^{\mathrm{Ad}}_{\ckG}(\varphi))=0$。
\item $\ckg^{I_\infty}=0$（注意 $\ckg$ 是 $V=\mathrm{Ad}$ 的空间）。
\end{enumerate}
以 $\Kl^{\theta^\vee}_{\ckG}(\varphi)$ 代替 $\Kl^{\mathrm{Ad}}_{\ckG}(\varphi)$ 时同样成立，且对 $\operatorname{char}(k)$ 无限制。
\end{proposition}

\begin{proof}
这些陈述是几何的，故在证明中忽略 Tate 扭转。由于伴随表示 $\ckg$ 自对偶，局部系统 $\Kl^{\mathrm{Ad}}_{\ckG}(\varphi)$ 也自对偶。固定这样的同构 $\Kl^{\mathrm{Ad}}_{\ckG}(\varphi)\xrightarrow{\sim}(\Kl^{\mathrm{Ad}}_{\ckG}(\varphi))^\vee$，得到同构
\begin{equation}
\begin{aligned}
H^0(\PP^1,j_{!*}\Kl^{\mathrm{Ad}}_{\ckG}(\varphi))&\cong H^2(\PP^1,j_{!*}\Kl^{\mathrm{Ad}}_{\ckG}(\varphi))^\vee,\\
H^0(\PP^1\setminus\{0,\infty\},\Kl^{\mathrm{Ad}}_{\ckG}(\varphi))&\cong H^2_c(\PP^1\setminus\{0,\infty\},\Kl^{\mathrm{Ad}}_{\ckG}(\varphi))^\vee.
\end{aligned}
\end{equation}
在 $D^b_c(\PP^1,\mathbb{Q}_\ell)$ 中有 distinguished triangle
\begin{equation}
j_!\Kl^{\mathrm{Ad}}_{\ckG}(\varphi)\to j_{!*}\Kl^{\mathrm{Ad}}_{\ckG}(\varphi)\to H^0i_0^*j_*\Kl^{\mathrm{Ad}}_{\ckG}(\varphi)\oplus H^0i_\infty^*j_*\Kl^{\mathrm{Ad}}_{\ckG}(\varphi)\to,
\end{equation}
它诱导长正合列
\begin{equation}
\begin{aligned}
0=H^0_c(\PP^1\setminus\{0,\infty\},\Kl^{\mathrm{Ad}}_{\ckG}(\varphi))&\to H^0(\PP^1,j_{!*}\Kl^{\mathrm{Ad}}_{\ckG}(\varphi))\to\ckg^{I_0}\oplus\ckg^{I_\infty}\xrightarrow{d}\\
&\xrightarrow{d}H^1_c(\PP^1\setminus\{0,\infty\},\Kl^{\mathrm{Ad}}_{\ckG}(\varphi))\to H^1_c(\PP^1,j_{!*}\Kl^{\mathrm{Ad}}_{\ckG}(\varphi))\to 0\to\\
&\to H^2_c(\PP^1\setminus\{0,\infty\},\Kl^{\mathrm{Ad}}_{\ckG}(\varphi))\to H^2_c(\PP^1,j_{!*}\Kl^{\mathrm{Ad}}_{\ckG}(\varphi))\to 0.
\end{aligned}
\end{equation}
由第二个正合列先得
\[
\begin{aligned}
H^2(\PP^1,j_{!*}\Kl^{\mathrm{Ad}}_{\ckG}(\varphi))&\cong H^2_c(\PP^1\setminus\{0,\infty\},\Kl^{\mathrm{Ad}}_{\ckG}(\varphi))\\
&\cong H^0(\PP^1\setminus\{0,\infty\},\Kl^{\mathrm{Ad}}_{\ckG}(\varphi))^\vee=0,
\end{aligned}
\]
最后一项的消没来自把引理 5.2 用于 $V=\ckg$。由对偶性，$H^0(\PP^1,j_{!*}\Kl^{\mathrm{Ad}}_{\ckG}(\varphi))=0$。由第一个正合列与 $H^0$ 的消没，连接同态 $d$ 是单射。

一方面，由定理 1(2)，$\dim\ckg^{I_0}=r(\ckG)$，因为 $I_0$ 通过主幂幺元作用在 $\ckg$ 上。另一方面，由定理 4，$\dim H^1_c(\PP^1\setminus\{0,\infty\},\Kl^{\mathrm{Ad}}_{\ckG}(\varphi))=-\chi_c(\PP^1\setminus\{0,\infty\},\Kl^{\mathrm{Ad}}_{\ckG}(\varphi))=r(\ckG)$（这里再次用了 $H^i_c(\PP^1\setminus\{0,\infty\},\Kl^{\mathrm{Ad}}_{\ckG}(\varphi))$ 对 $i=0,2$ 的消没）。因此 $d$ 必是同构。这蕴涵 $\ckg^{I_\infty}=0$。由第二个正合列，得到 $H^1(\PP^1,j_{!*}\Kl^{\mathrm{Ad}}_{\ckG}(\varphi))=0$。

对 $\Kl^{\theta^\vee}_{\ckG}(\varphi)$ 的陈述同样证明。由引理 10.1 证明中的等式 (10.2)，可以把引理 5.2 用于 $V_{\theta^\vee}$。
\end{proof}

\begin{remark}
沿用 Katz [27, §5.0]，称 $\PP^1\setminus\{0,\infty\}\otimes_k\bar k$ 上的 $\ckG$-局部系统 $L$ 为\emph{上同调刚性的}，若 $H^1(\PP^1,j_{!*}L^{\mathrm{Ad}})=0$。我们把 $H^1(\PP^1,j_{!*}L^{\mathrm{Ad}})$ 设想为 $L$ 的、在 $0$ 与 $\infty$ 周围形式去心圆盘上同构型固定的无穷小变形空间，尽管这种变形的概念尚未定义。命题 5.3 蕴涵 $\Kl_{\ckG}(\varphi)$ 是上同调刚性的，这为其物理刚性提供了证据。更精确的猜想见 §7。
\end{remark}

\subsection{一般计算方法}

以下计算中只考虑 $\Bun_{G(1,2)}$ 的中性分支。故可假设 $G$ 单连通。记 $\star\in\Bun_{G(1,2)}$ 为平凡丛对应的点。为计算层 $\Kl^V_{\ckG}(\varphi,\chi)$，把卷积图 (2.6) 限制到 $\star\times\PP^1\setminus\{0,\infty\}\subset\Bun_{G(1,2)}\times\PP^1\setminus\{0,\infty\}$：
\[
\begin{tikzcd}[column sep=large, row sep=large]
& \mathrm{GR} \arrow[dl, "\mathrm{pr}_1"'] \arrow[dr, "\mathrm{pr}_2"] & \\
\Bun_{G(1,2)} & & \star\times\PP^1\setminus\{0,\infty\}.
\end{tikzcd}
\]
由真基变换，$\Kl^V_{\ckG}(\varphi,\chi)=\mathrm{pr}_{2,!}(\mathrm{pr}_1^*A\otimes\mathrm{IC}_V)$。此外，$A$ 只支撑在大胞腔 $j:T\times I(1)/I(2)\subset\Bun_{G(1,2)}$ 上。记 $\mathrm{GR}^\circ\subset\mathrm{GR}$ 为大胞腔的原像，并写 $I(1)/I(2)\cong U/[U,U]\times U_{-\theta}$。按这个分解写 $\varphi=(\varphi_+,\varphi_0)$。

得到图
\[
\begin{tikzcd}[column sep=large, row sep=large]
\mathrm{GR}^\circ \arrow[r, hook] \arrow[dr, "\mathrm{pr}^\circ_2"'] \arrow[d, "{(f_T,f_0,f_+)}"'] & \mathrm{GR} \arrow[dl, "\mathrm{pr}_1"'] \arrow[dr, "\mathrm{pr}_2"] & \\
T\times U_{-\theta}\times U/[U,U] \arrow[r, hook] & \Bun_{G(1,2)} & \PP^1\setminus\{0,\infty\}
\end{tikzcd}
\]
由真基变换，
\begin{equation}
\Kl^V_{\ckG}(\varphi,\chi)=\mathrm{pr}^\circ_{2,!}\big(f_T^*\mathrm{Kum}_\chi\,f_0^*\mathrm{AS}_{\varphi_0}\otimes f_+^*\mathrm{AS}_{\varphi_+}\otimes\mathrm{IC}_V\big).
\end{equation}
故需描述子空间 $\mathrm{GR}^\circ$ 并计算映射 $f_T,f_0,f_+$。

首先记 $\mathrm{GR}^{\mathrm{triv}}\subset\mathrm{GR}$ 为遗忘映射 $\mathrm{GR}\to\Bun_G$ 下平凡 $G$-丛的原像。固定点 $x\in\PP^1\setminus\{0,\infty\}$，取 $t_x:=1-\tfrac tx$ 作为 $x$ 处的局部参数。于是 $t_x=1$ 对应 $t=0$，$t_x=\infty$ 对应 $t=\infty$。ind-概形 $\mathrm{Gr}^{\mathrm{triv}}_x$ 分类同构类对 $(E,\phi)$，其中 $E$ 是 $\PP^1$ 上的平凡 $G$-丛，$\phi:E|_{\PP^1\setminus\{x\}}\xrightarrow{\sim}G\times(\PP^1\setminus\{x\})$ 是同构。我们可以刚性化这个模问题，使 $\mathrm{Gr}^{\mathrm{triv}}_x$ 分类 $\PP^1\setminus\{x\}$ 上平凡 $G$-丛的、在 $\infty$ 处为恒等的自同构。故可等同
\begin{equation}
\mathrm{Gr}^{\mathrm{triv}}_x\xrightarrow{\ \sim\ }G[t_x^{-1}]_1:=\ker\big(G[t_x^{-1}]\xrightarrow{\mathrm{ev}(t_x=\infty)}G\big),
\end{equation}
其中 $G[t_x^{-1}]$ 是其 $R$-点为 $G(R[t_x^{-1}])$ 的 ind-概形。

\begin{remark}
有一族膨胀作用 $\mathbb{G}_m^{\mathrm{rot}}$ 作用在 $\PP^1\setminus\{0,\infty\}$ 上，它扩充到 $\mathrm{GR}$：$\lambda\in\mathbb{G}_m^{\mathrm{rot}}$ 把 $(x,E,\phi)$ 送到 $(\lambda x,\lambda^{-1,*}E,\lambda^{-1,*}\phi)$，并稳定 $\mathrm{GR}^{\mathrm{triv}}$。注意局部坐标 $t_x$ 在 $t$ 与 $x$ 的同时膨胀下不变。
\end{remark}

设 $\mathrm{Gr}^{\mathrm{triv}}\subset\mathrm{Gr}=G((\tau))/G[[\tau]]$ 是仿射 Grassmannian 的一个抽象拷贝，不涉及 $\PP^1\setminus\{0,\infty\}$ 上的任何点。膨胀作用连同 (5.9) 给出族 $\mathrm{GR}^{\mathrm{triv}}$ 在 $\PP^1\setminus\{0,\infty\}$ 上的平凡化：
\begin{equation}
\mathrm{GR}^{\mathrm{triv}}\cong(\PP^1\setminus\{0,\infty\})\times G[\tau^{-1}]_1\cong(\PP^1\setminus\{0,\infty\})\times\mathrm{Gr}^{\mathrm{triv}},\quad
(x,g(t_x^{-1}))\mapsto(x,g(\tau^{-1})).
\end{equation}
按定义，$\Bun_{G(1,2)}$ 中的大胞腔由 $T\times I(1)/I(2)$ 作用在平凡丛上得到。记 $\mathrm{ev}_0:G[t_x^{-1}]\to G[[t]]$ 与 $\mathrm{ev}_\infty:G[t_x^{-1}]\to G[[t^{-1}]]$ 为在 $0$ 与 $\infty$ 处的展开，$\mathrm{ev}_0,\mathrm{ev}_\infty:G[t_x^{-1}]\to G$ 为在 $t=0$ 与 $t=\infty$ 处的取值。

$G$-丛 $(E,\phi)\in\mathrm{Gr}_x$ 上的 $G(1,2)$-水平结构由平凡化 $\phi$ 得到。对 $g\in G[t_x^{-1}]$，复合 $E|_{\PP^1\setminus\{x\}}\xrightarrow{\phi}G\times\PP^1\setminus\{x\}\xrightarrow{g}G\times\PP^1\setminus\{x\}$ 扩充为同构 $E\to G\times\PP^1$。故 $g$ 定义了平凡丛上由 $\mathrm{ev}_0(g),\mathrm{ev}_\infty(g)$ 给出的水平结构。有交换图
\[
\begin{tikzcd}[column sep=large, row sep=large]
G[t_x^{-1}] \arrow[r, "{(\mathrm{ev}_0,\mathrm{ev}_\infty)}"] \arrow[dr, "{g\mapsto(\mathrm{ev}_0(g),\mathrm{ev}_\infty(g))}"'] & G\backslash\big(G[[t]]\times G[[t^{-1}]]\big) \arrow[d] \\
& \Bun_{G(1,2)}.
\end{tikzcd}
\]
这里 $G=\operatorname{Aut}(G\times\PP^1)$ 对角地作用在 $G[[t]]\times G[[t^{-1}]]$ 上。

于是我们发现 $g\in G[t_x^{-1}]$ 落在 $\mathrm{Gr}^\circ_x$ 中，当且仅当 $(\mathrm{ev}_0(g),\mathrm{ev}_\infty(g))\in G\backslash G(B^{\mathrm{opp}}\times U)$。利用等同 (5.9) 与 (5.11)，得到
\begin{equation}
\begin{aligned}
\mathrm{Gr}^\circ_x&\xrightarrow{\ \sim\ }\big\{g(t_x^{-1})\in G[t_x^{-1}]_1\ \big|\ \mathrm{ev}_0(g)\in UB^{\mathrm{opp}}\big\},\\
\mathrm{GR}^\circ&\xrightarrow{\ \sim\ }(\PP^1\setminus\{0,\infty\})\times\big\{g(\tau^{-1})\in G[\tau^{-1}]_1\ \big|\ g(1)\in UB^{\mathrm{opp}}\big\}
\end{aligned}
\end{equation}
包含在 $\mathrm{GR}^{\mathrm{triv}}=(\PP^1\setminus\{0,\infty\})\times G[\tau^{-1}]_1$ 中。

\begin{lemma}
对 (5.12) 参数化下的 $(x,g(\tau^{-1}))\in\mathrm{GR}^\circ$，写 $g(1)=ub^{\mathrm{opp}}$，其中 $u\in U$，$b^{\mathrm{opp}}\in B^{\mathrm{opp}}$。则
\[
\begin{aligned}
f_T(x,g)&=b^{\mathrm{opp}}\bmod U^{\mathrm{opp}}\in T;\qquad
f_+(x,g)=u^{-1}\bmod[U,U]\in U/[U,U];\\
f_0(x,g)&=x\,a_{-\theta}(g)\in U_{-\theta}\cong\ckg_{-\theta},
\end{aligned}
\]
其中 $a_{-\theta}:G[\tau^{-1}]\to\ckg_{-\theta}$ 把 $g$ 送到切向量 $\frac{dg(\tau^{-1})}{d(\tau^{-1})}\big|_{\tau^{-1}=0}\in\ckg$ 的 $\ckg_{-\theta}$-分量。
\end{lemma}

\begin{proof}
$f_T,f_+$ 的公式由 (5.12) 中的描述推出。按定义，$f_0(x,g)$ 由用局部参数 $t^{-1}$ 在 $t=\infty$ 处展开 $g(t_x^{-1})$ 并取 $t^{-1}$ 系数的 $\ckg_{-\theta}$-分量得到。注意
\[
\frac{dg(t_x^{-1})}{d(t^{-1})}\Big|_{t^{-1}=0}=x\,\frac{dg(t_x^{-1})}{d(t_x^{-1})}\Big|_{t_x^{-1}=0}\in\ckg.
\]
此外，在等同 (5.11) 下参数 $t_x$ 对应 $\tau$，故 $f_0(x,g)=x\,a_{-\theta}(g)$。引理得证。
\end{proof}

设 $\mathrm{AS}_{U,-\varphi_+}$ 为 $\mathrm{AS}_\psi$ 沿 $U\to U/[U,U]\xrightarrow{-\varphi_+}\mathbb{G}_a$ 的拉回。设 $\mathrm{Kum}_{B^{\mathrm{opp}},\chi}$ 为 $\mathrm{Kum}_\chi$ 沿 $B^{\mathrm{opp}}\to T$ 的拉回。设 $j_{UB^{\mathrm{opp}}}:UB^{\mathrm{opp}}\hookrightarrow G$ 为包含，并记
\[
J=J_{-\varphi_+,\chi}:=j_{UB^{\mathrm{opp}},!}\big(\mathrm{AS}_{U,-\varphi_+}\boxtimes\mathrm{Kum}_{B^{\mathrm{opp}},\chi}\big)\in D^b_c(G,\mathbb{Q}_\ell).
\]

\begin{remark}
由 [4]，或本文引理 2.3 的论证，$J$ 有 clean 性：
\[
j_{UB^{\mathrm{opp}},!}\big(\mathrm{AS}_{U,-\varphi_+}\boxtimes\mathrm{Kum}_{B^{\mathrm{opp}},\chi}\big)\xrightarrow{\ \sim\ }j_{UB^{\mathrm{opp}},*}\big(\mathrm{AS}_{U,-\varphi_+}\boxtimes\mathrm{Kum}_{B^{\mathrm{opp}},\chi}\big).
\]
我们可以把 $J=J_{-\varphi_+,\chi}\in D^b_c(G,\mathbb{Q}_\ell)$ 视作定义 2.4 与注 2.5 中自守层 $A_{\varphi,\chi}$ 的有限域类比。\zhnote{原文此处引作 ``Definition 2.2''，但该编号在 §2 中已被注 2.2 占用，所指显然是定义 2.4 的自守层；按 2.4 译。[来源：原文 p.30 与 §2.2 的编号顺序]}
\end{remark}

有了这些记号并利用等同 (5.11)，公式 (5.8) 变为
\begin{equation}
\Kl^V_{\ckG}(\varphi,\chi)\xrightarrow{\ \sim\ }\pi^{\mathrm{triv}}_!\big(\mathrm{IC}^{\mathrm{triv}}_V\otimes f_0^*\mathrm{AS}_{\varphi_0}\otimes\mathrm{ev}^*_{\tau=1}J\big),
\end{equation}
其中 $\pi^{\mathrm{triv}}:\mathrm{GR}^{\mathrm{triv}}=\PP^1\setminus\{0,\infty\}\times G[\tau^{-1}]_1\to\PP^1\setminus\{0,\infty\}$ 是投影。以下常把 $\mathrm{ev}_{\tau=1}$ 简写作 $\mathrm{ev}$。

本节最终目标是计算 $V=V_{\theta^\vee}$ 与 $\ckg$ 时 $\Kl^V_{\ckG}(\varphi,\chi)$ 的 Euler 特征。这里先对一般 $V$ 作几步化归。由 (5.15)，需计算
\begin{equation}
\chi_c(\PP^1\setminus\{0,\infty\},\Kl^V_{\ckG}(\varphi,\chi))=\chi_c\big(\PP^1\setminus\{0,\infty\}\times\mathrm{Gr}^{\mathrm{triv}},\mathrm{IC}^{\mathrm{triv}}_V\otimes f_0^*\mathrm{AS}_{\varphi_0}\otimes\mathrm{ev}^*J\big).
\end{equation}
按 $a_{-\theta}$ 是否消没，把 $\mathrm{Gr}^{\mathrm{triv}}\cong G[\tau^{-1}]_1$ 分解为 $\mathrm{Gr}^{\mathrm{triv},a\neq 0}$ 与 $\mathrm{Gr}^{\mathrm{triv},a=0}$。在 $\PP^1\setminus\{0,\infty\}\times\mathrm{Gr}^{\mathrm{triv},a=0}$ 上，(5.16) 中的复形沿 $\PP^1\setminus\{0,\infty\}$ 常值，故 Euler 特征为 $0$。

另一方面，有变量替换同构
\[
\PP^1\setminus\{0,\infty\}\times\mathrm{Gr}^{\mathrm{triv},a\neq 0}\ni(x,g)\mapsto(xa_{-\theta}(g),g)\in\mathbb{G}_m\times\mathrm{Gr}^{\mathrm{triv},a\neq 0}.
\]
在此同构下，$f_0=xa_{-\theta}$ 变成到 $\mathbb{G}_m$-因子的投影，于是可应用 Künneth 公式。汇总这些步骤，得到
\begin{align}
\chi_c(\PP^1\setminus\{0,\infty\},\Kl^V_{\ckG}(\varphi,\chi))
&=\chi_c\big(\PP^1\setminus\{0,\infty\}\times\mathrm{Gr}^{\mathrm{triv},a\neq 0},\mathrm{IC}^{\mathrm{triv}}_V\otimes f_0^*\mathrm{AS}_{\varphi_0}\otimes\mathrm{ev}^*J\big)\notag\\
&\qquad=\chi_c(\mathbb{G}_m,\mathrm{AS}_{\varphi_0})\,\chi_c(\mathrm{Gr}^{\mathrm{triv},a\neq 0},\mathrm{IC}^{\mathrm{triv}}_V\otimes\mathrm{ev}^*J)\notag\\
&=-\chi_c(\mathrm{Gr}^{\mathrm{triv},a\neq 0},\mathrm{IC}^{\mathrm{triv}}_V\otimes\mathrm{ev}^*J)\notag\\
&=-\chi_c(\mathrm{Gr}^{\mathrm{triv}}\setminus\star,\mathrm{IC}^{\mathrm{triv}}_V\otimes\mathrm{ev}^*J)\notag\\
&\qquad+\chi_c(\mathrm{Gr}^{\mathrm{triv},a=0}\setminus\star,\mathrm{IC}^{\mathrm{triv}}_V\otimes\mathrm{ev}^*J).
\end{align}
最后一个等式是因为基点 $\star\in\mathrm{Gr}$ 属于 $\mathrm{Gr}^{\mathrm{triv},a=0}$。

\begin{lemma}
$\chi_c(\mathrm{Gr}^{\mathrm{triv}}\setminus\star,\mathrm{IC}^{\mathrm{triv}}_V\otimes\mathrm{ev}^*J)=0$。
\end{lemma}

\begin{proof}
由于 $\mathrm{IC}_V$ 沿层 $\mathrm{Gr}_\lambda$ 常值，只需证明对支配余权 $\lambda\neq 0$ 有 $\chi_c(\mathrm{Gr}^{\mathrm{triv}}_\lambda,\mathrm{ev}^*J)=0$。

记 $D^b_c(G,\mathbb{Q}_\ell)$ 上的卷积积为：对 $K_1,K_2\in D^b_c(G,\mathbb{Q}_\ell)$，
\[
K_1*_G K_2:=m_!(K_1\boxtimes K_2),
\]
其中 $m:G\times G\to G$ 是乘法。令 $K_\lambda:=\mathrm{ev}_!\mathbb{Q}_{\ell,\mathrm{Gr}^{\mathrm{triv}}_\lambda}\in D^b_c(G)$。则 $R\Gamma_c(\mathrm{Gr}^{\mathrm{triv}}_\lambda,\mathrm{ev}^*J)$ 是 $\mathrm{AS}_{U,\varphi_+}*_G K_\lambda*_G\mathrm{Kum}_{B^{\mathrm{opp}},\chi^{-1}}$ 在 $e\in G$ 处的茎。由于 $\mathrm{ev}=\mathrm{ev}_{\tau=1}$ 是 $G$-等变的（在共轭下），$K_\lambda$ 带有自然的 $G$-等变结构。故
\[
K_\lambda*_G\mathrm{Kum}_{B^{\mathrm{opp}},\chi^{-1}}\cong\mathrm{Kum}_{B^{\mathrm{opp}},\chi^{-1}}*_G K_\lambda.
\]
特别地，在每个 Bruhat 层 $B^{\mathrm{opp}}wB^{\mathrm{opp}}$ 上，$K_\lambda*_G\mathrm{Kum}_{B^{\mathrm{opp}},\chi^{-1}}$ 形如 $S_w\otimes L_w$，其中 $L_w$ 是 $U^{\mathrm{opp}}wB^{\mathrm{opp}}$ 上的局部系统，而 $S_w$（$\mathbb{Q}_\ell$-向量空间复形）是 $K_\lambda*_G\mathrm{Kum}_{B^{\mathrm{opp}},\chi^{-1}}$ 在 $\dot w\in N_G(T)$（$w\in W$ 的代表元）处的茎。

因此 $\mathrm{AS}_{U,\varphi_+}*_G K_\lambda*_G\mathrm{Kum}_{B^{\mathrm{opp}},\chi^{-1}}$ 是诸 $S_w\otimes(\mathrm{AS}_{U,\varphi_+}*_G L_w)$（$w\in W$）的逐次扩张。为证明 $\mathrm{AS}_{U,\varphi_+}*_G K_\lambda*_G\mathrm{Kum}_{B^{\mathrm{opp}},\chi^{-1}}$ 的茎的 Euler 特征消没，只需证明对所有 $w$ 有 $\chi(S_w)=0$。

按定义，
\begin{equation}
\begin{split}
S_w&=\big(K_\lambda*_G\mathrm{Kum}_{B^{\mathrm{opp}},\chi^{-1}}\big)_{\dot w}=R\Gamma_c(wB^{\mathrm{opp}},K_\lambda\otimes\mathrm{Kum}_{B^{\mathrm{opp}},\chi})\\
&=R\Gamma_c\big(\mathrm{Gr}^{\mathrm{triv}}_\lambda\cap\mathrm{ev}^{-1}(wB^{\mathrm{opp}}),\mathrm{ev}^*\mathrm{Kum}_{B^{\mathrm{opp}},\chi}\big).
\end{split}
\end{equation}
$T$ 通过共轭作用在 $\mathrm{Gr}^{\mathrm{triv}}=G[\tau^{-1}]_1$ 上，并保持 $\mathrm{Gr}^{\mathrm{triv}}_\lambda\cap\mathrm{ev}^{-1}(wB^{\mathrm{opp}})$。$\mathrm{Gr}$ 上唯一的 $T$-不动点是 $\tau^\mu$（$\mu\in X_*(T)$），它们不属于任何 $\mathrm{Gr}^{\mathrm{triv}}_\lambda$（$\lambda\neq 0$）。此外，局部系统 $\mathrm{ev}^*\mathrm{Kum}_{B^{\mathrm{opp}},\chi}$ 在 $T$-共轭下是单值的：存在 $m\ge 1$（与 $p$ 互素）使 $\mathrm{ev}^*\mathrm{Kum}_{B^{\mathrm{opp}},\chi}$ 在 $T$-共轭的 $m$ 次幂下等变。由于 $T$-共轭的 $m$ 次幂在 $\mathrm{Gr}^{\mathrm{triv}}_\lambda\cap\mathrm{ev}^{-1}(wB^{\mathrm{opp}})$ 上仍无不动点，(5.19) 中的 Euler 特征为零。引理得证。
\end{proof}

\begin{corollary}
设 $X_*(T)^+\subset X_*(T)$ 为支配余权，$V(\lambda)\subset V$ 为权 $\lambda\in X_*(T)$ 的权空间。则
\[
\chi_c(\PP^1\setminus\{0,\infty\},\Kl^V_{\ckG}(\varphi,\chi))=\sum_{\lambda\in X_*(T)^+,\ \lambda\neq 0}\dim V(\lambda)\ \chi_c\big(\mathrm{Gr}^{\mathrm{triv},a=0}_\lambda,\mathrm{ev}^*J\big).
\]
\end{corollary}

\begin{proof}
由 [29, Theorem 6.1]，$\mathrm{IC}_V$ 沿 $\mathrm{Gr}_\lambda$ 的茎的 Euler 特征为 $\dim V(\lambda)$。因此
\[
\chi_c\big(\mathrm{Gr}^{\mathrm{triv},a=0}\setminus\star,\mathrm{IC}^{\mathrm{triv}}_V\otimes\mathrm{ev}^*J\big)=\sum_{\lambda\in X_*(T)^+,\ \lambda\neq 0}\dim V(\lambda)\ \chi_c\big(\mathrm{Gr}^{\mathrm{triv},a=0}_\lambda,\mathrm{ev}^*J\big).
\]
把它与 (5.17) 及引理 5.18 结合即得所求等式。
\end{proof}
